\documentclass[11pt]{article}
\usepackage[letterpaper,margin=0.82in,headheight=22pt,footskip=28pt]{geometry}
\usepackage[T1]{fontenc}
\usepackage{mathptmx}
\usepackage{microtype}
\usepackage{amsmath,amssymb,mathtools,bm}
\usepackage{graphicx}
\usepackage{booktabs,tabularx,longtable,array,multirow}
\usepackage{enumitem}
\usepackage[dvipsnames]{xcolor}
\usepackage{fancyhdr}
\usepackage{titlesec}
\usepackage{caption,subcaption}
\usepackage[numbers,sort&compress]{natbib}
\usepackage{hyperref}
\usepackage{url}
\usepackage{float}
\usepackage{placeins}
\usepackage{setspace}
\hypersetup{colorlinks=true,linkcolor=NavyBlue,citecolor=NavyBlue,urlcolor=MidnightBlue,pdftitle={Persistent Landscapes, Hodge Curl, and Schrödinger-Bridge Reinforcement Learning for Optimal Cell-Fate Control},pdfauthor={Devanik Debnath}}
\definecolor{Ink}{HTML}{142235}
\definecolor{Accent}{HTML}{236A73}
\definecolor{Muted}{HTML}{607184}
\definecolor{Light}{HTML}{EEF4F5}
\titleformat{\section}{\Large\bfseries\color{Ink}}{\thesection}{0.65em}{}
\titleformat{\subsection}{\large\bfseries\color{Accent}}{\thesubsection}{0.65em}{}
\titleformat{\subsubsection}{\normalsize\bfseries\color{Ink}}{\thesubsubsection}{0.6em}{}
\titlespacing*{\section}{0pt}{2.0ex plus .4ex minus .2ex}{1.0ex}
\titlespacing*{\subsection}{0pt}{1.5ex plus .3ex minus .2ex}{0.6ex}
\captionsetup{font=small,labelfont={bf,color=Ink},textfont=small,skip=6pt}
\setlist{itemsep=3pt,topsep=4pt,leftmargin=1.5em}
\setstretch{1.06}
\pagestyle{fancy}
\fancyhf{}
\fancyhead[L]{\small\color{Muted}Persistent Landscapes and Cell-Fate Control}
\fancyhead[R]{\small\color{Muted}Theory--simulation preprint}
\fancyfoot[C]{\small\thepage}
\renewcommand{\headrulewidth}{0.4pt}
\renewcommand{\headrule}{\hbox to\headwidth{\color{Accent}\leaders\hrule height \headrulewidth\hfill}}
\newcommand{\E}{\mathbb{E}}
\newcommand{\R}{\mathbb{R}}
\newcommand{\Prob}{\mathbb{P}}
\newcommand{\D}{\mathrm{d}}
\newcommand{\KL}{D_{\mathrm{KL}}}
\newcommand{\argmin}{\mathop{\mathrm{arg\,min}}}
\newcommand{\TV}{\mathrm{TV}}
\newcommand{\curl}{\operatorname{curl}}
\newcommand{\grad}{\nabla}
\newcommand{\tr}{\mathrm{tr}}
\newcommand{\1}{\mathbf{1}}
\newcommand{\Hzero}{H_0}
\newcommand{\Hone}{H_1}
\newcommand{\claim}[1]{\medskip\noindent\colorbox{Light}{\parbox{0.96\linewidth}{\textbf{Scope note.} #1}}\medskip}
\newcommand{\fig}[4]{\begin{figure}[tbp]\centering\includegraphics[width=#1\linewidth]{figures/#2.pdf}\caption{#3}\label{#4}\end{figure}}
\begin{document}
\begin{titlepage}
\thispagestyle{empty}
\vspace*{0.24in}
\noindent{\color{Accent}\rule{\linewidth}{2pt}}\par\vspace{0.10in}
{\fontsize{19}{23}\selectfont\bfseries\color{Ink}\parbox{\dimexpr\linewidth-22pt\relax}{Persistent Landscapes, Hodge Curl, and\\[0.10in] Schrödinger-Bridge\\[0.10in] Reinforcement Learning for Optimal Cell-Fate Control}\par}
\vspace{0.10in}
{\Large\color{Accent}Topological Obstructions to Cellular Reprogramming\par}
\vspace{0.18in}
{\large Devanik Debnath (Devanik21)\par}
\vspace{0.05in}
{\normalsize Department of Electronics and Communication Engineering\\National Institute of Technology Agartala, India\par}
\vspace{0.12in}
{\small
\noindent\textbf{Manuscript type:} Theory-plus-simulation research preprint\par
\noindent\textbf{Version:} v2026.10.24.1\par
\noindent\textbf{Date:} 24 October 2026\par
\noindent\textbf{Status:} Synthetic computational study; no empirical single-cell data fitted\par
\noindent\textbf{Scope:} Mathematical analysis and falsifiable hypotheses; not a validated cellular protocol\par
}
\vspace{0.16in}
\noindent{\color{Accent}\rule{\dimexpr\linewidth-16pt\relax}{0.8pt}}\par\vspace{0.12in}
\textbf{Abstract}\par\vspace{0.04in}\begingroup\fontsize{9.3}{11.0}\selectfont
A cell-fate landscape is a useful summary, but it is not automatically a thermodynamic potential or a control-cost surface. This paper tests a topology-aware workflow linking persistent homology, stochastic action, Hodge decomposition, Schrödinger bridges, and reinforcement learning. In a one-dimensional quartic gradient model, the analytic barrier is 0.25 and the matched H-zero persistence estimate is 0.249998, conditional on correct minimum--saddle pairing. Under observation noise, persistence calibration and pairing degrade. An analytic circle example with first Betti number one shows that constant control can minimize effort: a matched-endpoint pulse costs 3.33 times more, falsifying the unconditional claim that a loop forces pulsing. In 400 synthetic rollouts per policy, evenly spaced pulses succeed in 23.25\% of episodes, while the tabular RL policy achieves 7.75\% and constant dosing 0\% in the chosen regime. RL therefore loses to a simple heuristic. All results are synthetic; no single-cell data were analyzed, no wet-lab experiments were run, and no biological efficacy is claimed.
\par\vspace{0.05in}
\noindent\textbf{Keywords:} cell-fate control; persistent homology; Freidlin--Wentzell action; Hodge decomposition; Schrödinger bridge; stochastic optimal control; reinforcement learning; synthetic validation
\par\vspace{0.04in}\noindent{\small\color{Muted}\textit{Scope: synthetic mathematical benchmark only; not a validated cellular protocol.}}\endgroup
\end{titlepage}
\pagenumbering{roman}
\setcounter{tocdepth}{1}
\tableofcontents
\clearpage
\pagenumbering{arabic}

\section{Research question and contribution}
Cell identity is not directly a point in a one-dimensional potential. It is a high-dimensional and partially observed state shaped by transcriptional regulation, chromatin accessibility, signaling, metabolic context, cell-cycle state, and stochastic molecular events. The landscape metaphor is useful because it makes stability, transition routes, and competing attractors discussable in a common geometric language. It is also dangerous: a plotted scalar field may be an inferred summary rather than a thermodynamic potential, and an attractive visual metaphor can be mistaken for a quantitative law.

This paper asks whether topological summaries can impose useful constraints on minimum-effort control between cell-state distributions. The motivation is practical. Reprogramming protocols typically involve time-varying transcription-factor expression, media changes, or other interventions whose sequence, amplitude, and duration matter. Trial-and-error optimization is expensive. A data-driven controller might search a much larger schedule space, but the optimization becomes fragile if the state representation is noisy, the dynamical model is misidentified, or the objective ignores hidden paths.

The project begins with three connected claims: (i) basin barriers are recoverable as zero-dimensional persistence when a scalar gradient potential and the relevant minimum--saddle pairing are known; (ii) rotational drift and topological cycles may make temporal schedules more useful than constant inputs; and (iii) reinforcement learning (RL), regularized by stochastic path-control objectives, may discover schedules that exploit such structure. Each claim has a different epistemic status, so we do not bundle them as a single theorem.

The first part is a conditional mathematical identity, not a universal equivalence between persistence and stochastic control cost. The second is a conjecture that can fail: nonzero first homology alone does not require pulsed input. The third is an empirical question for a specified simulator. In our initial test, the RL policy does not outperform an evenly spaced pulse heuristic. This negative result changes the paper's central contribution from an asserted topology-to-protocol law to an audit of where such a law would need additional assumptions.

\subsection{Contributions and non-contributions}
The paper provides four concrete outputs. First, it specifies the assumptions under which a potential barrier and an H\textsubscript{0} persistence value coincide. Second, it gives an analytic circle counterexample to an unconditional loop-implies-pulsing statement. Third, it implements a deterministic, reproducible synthetic benchmark covering numerical action convergence, persistence under noise, neural drift fitting, and matched-budget protocol comparisons. Fourth, it records both favorable and unfavorable outputs and turns them into a falsifiable plan for eventual analysis of a public reprogramming time course.

It does not claim a new general theorem of cellular reprogramming; it does not claim that Hodge curl creates a particular biological control barrier; it does not report learning from real single-cell data; it does not claim RL superiority; and it does not propose a wet-lab dosing schedule. The principal methodological novelty claim is modest: an integrated, fail-explicit workflow combining conditional landscape persistence, a Hodge-aware dynamical interpretation, stochastic bridge objectives, and protocol benchmarks with explicit counterexamples. A systematic priority search could still reveal closer integrated approaches.

\begin{table}[tbp]\centering\small
\caption{Evidence taxonomy for the paper's main statements.}
\label{tab:epistemic}
\begin{tabularx}{\linewidth}{@{}p{0.13\linewidth}p{0.36\linewidth}X@{}}\toprule
Class & Meaning here & What the evidence supports \\\midrule
P1 & Mathematical identity under stated assumptions & Quartic barrier--persistence match; minimum-energy control on the circle \\
P2 & Variational or stochastic-control statement & FW action relation for the specified gradient SDE; bridge objective definition \\
P3 & Statistical or robustness claim & Synthetic noise sweeps and policy estimates within the tested parameter grid only \\
P4 & Cross-domain mapping & Using persistence and Hodge features as summaries for control; requires empirical validation \\
P5 & Intuition or conjecture & Cycles may indicate schedule sensitivity; pulse-count lower bound remains unproved and is not supported generally \\\bottomrule\end{tabularx}
\end{table}

\section{Prior work and novelty boundary}
\subsection{Landscape models of cell fate}
Classical developmental biology introduced the landscape metaphor as a way of visualizing constrained developmental trajectories. Modern quantitative treatments derive or estimate a landscape from dynamical models and use basin geometry to discuss stability and transitions. Work by Zhou and colleagues quantified landscape topography and paths for differentiation, reprogramming, and transdifferentiation, reporting a relation between barrier heights and transition rates in a specified gene-network model \citep{zhou2013waddington}. That history rules out a novelty claim for the broad association between cell-fate landscape barriers and fate-transition difficulty.

The relevant gap is narrower: an inferred scalar landscape and the actual stochastic control problem may differ in metric, noise covariance, hidden variables, and non-gradient flux. A persistence interval is a property of a filtered scalar function. A Freidlin--Wentzell quasipotential is a dynamical quantity defined relative to a drift and noise covariance. They coincide in special cases, but neither object can silently substitute for the other. This distinction is central to the present audit.

\subsection{Persistent homology and uncertainty}
Persistent homology records how connected components, loops, and higher-dimensional features appear and disappear as a filtration threshold changes. The foundational theory and stability results give ways to compare persistence diagrams under controlled perturbations \citep{edelsbrunner2002persistence,cohensteiner2007stability}. For a Morse potential, the elder-rule pairing in H\textsubscript{0} can associate a local minimum with the saddle level at which its component joins an older component. The interval length is then a minimum--saddle height difference. This is not automatically the barrier from an arbitrary source state to an arbitrary target basin: the pairing must describe that transition, and a scalar potential must exist.

Data-derived landscapes add estimation error. Smoothing, feature selection, dimensionality reduction, and sampling density affect critical points and their pairings. Persistence stability can control changes in the diagram under a uniform perturbation of the scalar function, but it does not guarantee that the learned function is the correct biological potential or that its minimum--saddle pairing corresponds to the relevant dynamic escape route. The noise experiment in this paper demonstrates a particularly blunt failure mode: with increasing observation noise, estimated persistence becomes biased upward and the target pairing is rarely identified by the simple grid-based procedure.

\subsection{Stochastic control and large deviations}
For small-noise diffusions, large-deviation theory assigns an action to paths and gives exponential-scale predictions for rare transitions. Kramers-type escape approximations relate transition rates to barrier structure under suitable assumptions \citep{kramers1940brownian,freidlin2012random}. The gradient-diffusion setting is especially interpretable: for an isotropic-noise SDE with drift equal to minus the gradient of a potential, the quasipotential cost to cross a saddle agrees with the potential increase in the usual noise normalization. In non-gradient systems, the quasipotential solves a dynamical variational problem and generally cannot be reconstructed from a convenient scalar projection alone.

\subsection{Hodge structure, optimal transport, and RL}
Hodge decomposition separates gradient-like, divergence-free, and harmonic components relative to a chosen domain, metric, boundary convention, and discretization. Related combinatorial formulations have been applied to ranking and cyclic inconsistency \citep{jiang2011hodge}. The decomposition does not equate every curl component to a nontrivial global homology class: local rotational drift can exist on a simply connected domain. Conversely, a space with nontrivial first Betti number does not imply that a particular control objective favors pulses.

Optimal transport methods infer population-level correspondences between time-stamped single-cell distributions. Schrödinger bridges formulate a stochastic analogue, selecting a path measure close in relative entropy to a reference process while meeting endpoint constraints. Contemporary work explicitly studies Schrödinger-problem formulations for cell-fate trajectory inference, including mechanistic bursty gene-expression references \citep{fournie2026schrodinger,deciphering2025cellfate}. Deep RL has also been proposed for cellular reprogramming in Boolean and probabilistic Boolean gene-regulatory networks: pbn-STAC is an especially close neighbor \citep{zarzycki2025pbnstac}. The 2026 GATTACA framework further combines graph neural networks with deep RL to control Boolean biological networks at scale \citep{mizera2026gattaca}. These systems establish that RL-based network control is not itself a novelty claim. The present study differs in its synthetic continuous-state emphasis and its explicit counterexample-driven audit of proposed links between persistence, non-gradient dynamics, and schedules; it does not show superiority over these prior methods.

\begin{table}[tbp]\centering\small
\caption{Closest literatures and the claim they do—and do not—license.}
\label{tab:related}
\begin{tabularx}{\linewidth}{@{}p{0.23\linewidth}p{0.34\linewidth}X@{}}\toprule
Literature & Established contribution & Boundary for this paper \\\midrule
Waddington / landscape models & Relate basin/barrier geometry to cell-fate stability and routes & Does not make every inferred landscape a thermodynamic potential \\
Persistent homology & Describes multiscale connectivity and has stability theorems & Does not identify dynamical escape cost without a model-to-function correspondence \\
Freidlin--Wentzell / Kramers & Action-based rare-event and escape-rate theory & Gradient barrier formulas depend on drift/noise assumptions \\
Hodge decomposition & Splits discrete or continuous fields into components under defined geometry & Curl is not synonymous with global H\textsubscript{1}; result depends on the domain and boundaries \\
Schrödinger bridge / OT & Matches distributions through low-relative-entropy stochastic paths & Endpoint transport alone does not establish causal cell trajectories \\
pbn-STAC and other RL approaches & Learn reprogramming interventions in network models & Our toy benchmark is not a superiority result over those methods \\
\bottomrule\end{tabularx}
\end{table}

\subsection{Research gap and positioning}
We find substantive overlap in each ingredient and related efforts combining cell-fate inference, stochastic transport, landscape estimation, and RL. The paper therefore does not claim the first use of topology, RL, or Schrödinger bridges in cellular dynamics. The narrower question is whether a carefully separated set of topological and dynamical diagnostics can predict when a schedule advantage is plausible and when that prediction should be rejected. This work is a first reproducible stress test, not a validated universal solution. Search coverage included primary literature, review and methods articles, DOI/PubMed/PMC records, the pbn-STAC code repository, and GEO accessions. It was a focused search, not an exhaustive patent or all-databases priority review.

\section{State space, controls, and assumptions}
Let $X_t\in\mathcal{X}\subseteq\R^d$ represent a latent cell-state embedding, not the full molecular state. Observations $Y_t$ are produced by an observation map $h$ with measurement noise. A control $u_t\in\mathcal{U}$ represents a coarse-grained intervention such as a transcription-factor induction level or media perturbation; the present simulation uses a scalar abstract control and does not map it to a biological dose.

We define the controlled diffusion
\begin{equation}
\D X_t = \bigl(b(X_t)+B(X_t)u_t\bigr)\D t + \sqrt{2\varepsilon}\,\sigma(X_t)\D W_t,
\label{eq:sde}
\end{equation}
where $b$ is the uncontrolled drift, $B$ maps control into state velocity, $\sigma$ determines noise geometry, and $\varepsilon$ controls noise scale. All are unknown in a real application and would need uncertainty estimates. In our simulator the state is one-dimensional and synthetic. The trained neural estimator approximates a known synthetic drift; it is not trained on experimental transcriptomics.

A protocol is a measurable control schedule satisfying amplitude, total exposure, and duration constraints. We separate a nominal dose budget $D(u)=\int_0^T\|u_t\|_1\,\D t$ from a quadratic effort $E(u)=\tfrac12\int_0^T u_t^\top R u_t\,\D t$. These quantities are not interchangeable. A matched L\textsubscript{1} dose does not guarantee matched biological toxicity, energy, or off-target effect. The synthetic benchmark matches nominal scalar exposure and reports the quadratic energy only in the analytic rotor example.

The endpoint is a target set $\mathcal{B}$ representing a desired state region. For an individual trajectory, success means $X_T\in\mathcal{B}$; a population objective may instead require a terminal distribution $\mu_T$ to match a reference $\nu_1$. The distinction matters because the latter values heterogeneity and can be more faithful to population reprogramming than optimizing a mean trajectory. No state region in the synthetic model is labelled “pluripotent”; it is a target set in an abstract bistable system.

\begin{table}[tbp]\centering\small
\caption{Notation and dimensional conventions.}
\label{tab:notation}
\begin{tabularx}{\linewidth}{@{}p{0.18\linewidth}p{0.30\linewidth}X@{}}\toprule
Symbol & Definition & Interpretation \\\midrule
$X_t$ & Latent state in $\mathcal{X}$ & Reduced expression-state coordinate or vector \\
$Y_t$ & Observation of $X_t$ & Noisy measured representation \\
$b(x)$ & Uncontrolled drift & State velocity before intervention \\
$B(x)u$ & Control contribution & Abstract intervention action \\
$\sigma(x)$, $\varepsilon$ & Diffusion geometry, noise scale & Must be estimated/calibrated \\
$V(x)$ & Scalar potential, when it exists & Not assumed for arbitrary drift \\
$\Delta V$ & $V(x_s)-V(x_m)$ & Potential rise from minimum to saddle \\
$\mathrm{Pers}_0$ & H\textsubscript{0} interval length & Death threshold minus birth threshold \\
$S(\phi)$ & FW action of path $\phi$ & Path cost under specified drift/noise \\
$\mu_t,\nu_t$ & Population probability measures & Marginal distributions at time $t$ \\
$D(u), E(u)$ & L\textsubscript{1} dose, quadratic effort & Distinct constraints/objectives \\
$\beta$ & Synthetic drift-shape parameter & Dimensionless simulator parameter \\
\bottomrule\end{tabularx}
\end{table}

\subsection{Assumptions that must be exposed}
The simple barrier identity assumes a smooth enough scalar potential with a nondegenerate local minimum and relevant index-one saddle, sublevel-set filtration of precisely that potential, and correct persistence pairing. The simple FW interpretation further assumes an isotropic-noise gradient diffusion with a stated normalization and a transition governed by the relevant saddle. The controller comparison assumes common initial-state distributions, horizon, action budget, and success definition. The neural model assumes a specified state representation and training sampling process. None of these assumptions is guaranteed by snapshot single-cell data.

\section{Barrier height, escape action, and H\textsubscript{0} persistence}
\subsection{Gradient setting}
Consider the small-noise diffusion
\begin{equation}
\D X_t=-\grad V(X_t)\D t+\sqrt{2\varepsilon}\,\D W_t.
\label{eq:gradientSDE}
\end{equation}
For an absolutely continuous path $\phi:[0,T]\to\mathcal{X}$, the Freidlin--Wentzell action under this normalization is
\begin{equation}
S_T(\phi)=\frac14\int_0^T\|\dot\phi(t)+\grad V(\phi(t))\|^2\,\D t.
\label{eq:action}
\end{equation}
For a path climbing from a minimum $x_m$ to a saddle $x_s$, the identity
\begin{align}
\frac14\|\dot\phi+\grad V\|^2
&=\frac14\|\dot\phi-\grad V\|^2+\dot\phi\cdot\grad V
\end{align}
shows that its action is bounded below by $V(x_s)-V(x_m)$; the lower bound is approached along the time-reversed deterministic descent trajectory when such a heteroclinic minimizer exists. Thus the minimum action equals $\Delta V$ in this idealized regime. The physical transition-rate exponent scales as $\Delta V/\varepsilon$ under Eq.~\eqref{eq:gradientSDE}; prefactors and finite-noise corrections are not encoded by the barrier alone.

\subsection{Persistence pairing}
Given the sublevel filtration $V^{-1}((-\infty,a])$, connected components are born at local minima and merge at saddle levels. If a particular minimum is paired to the index-one saddle separating it from the older component according to the elder rule, its H\textsubscript{0} persistence is
\begin{equation}
\mathrm{Pers}_0(x_m,x_s)=V(x_s)-V(x_m)=\Delta V.
\label{eq:persistence}
\end{equation}
Equation~\eqref{eq:persistence} is an equality of two height differences taken from the same scalar function. It is not a theorem that a persistence interval computed from an arbitrary learned score equals the minimal control energy, a biological energy, or the quasipotential barrier for non-gradient dynamics. For anisotropic noise, state-dependent diffusion, non-gradient drift, or incorrect source-saddle pairing, the relevant quasipotential is defined by the controlled action and may differ from raw sublevel persistence.

\fig{0.88}{fig01_landscape}{Synthetic quartic landscape $V(x)=(x^2-1)^2/4$. The minima and central saddle define a potential barrier of 0.25 in this example. The visualization is a synthetic mathematical object, not a fitted cell-state landscape.}{fig:landscape}

\fig{0.82}{fig02_barrier_persistence}{Conditional match between the analytic potential barrier and the grid-estimated H\textsubscript{0} persistence for the quartic example. Numerical agreement is a unit test of the specified scalar landscape and pairing; it does not validate a general biological mapping.}{fig:barrier}

\subsection{Numerical action check}
The verification script discretizes candidate paths and evaluates the action integral at successively refined resolutions. The absolute error declines from approximately $8.02\times10^{-5}$ on the coarsest grid to $1.48\times10^{-6}$ on the finest grid in the saved convergence test. This is evidence that the implemented quadrature converges toward the analytic result in this toy case. It is not an independent test of the gradient assumption or evidence that a real cell-state transition follows that path.

\begin{table}[tbp]\centering\small
\caption{Analytic and numerical checks for the quartic toy potential.}
\label{tab:barrierresults}
\begin{tabularx}{\linewidth}{@{}Xrr@{}}\toprule
Quantity & Value & Status \\\midrule
Potential minimum $V(-1)$ & 0.000000 & Analytic \\
Central saddle $V(0)$ & 0.250000 & Analytic \\
Analytic barrier $\Delta V$ & 0.250000 & Analytic \\
Grid persistence estimate & 0.249998 & Numerically recovered \\
Absolute barrier--persistence mismatch & $1.69\times10^{-6}$ & Synthetic unit check \\
Coarse-grid FW-action absolute error & $8.02\times10^{-5}$ & Numerical convergence \\
Fine-grid FW-action absolute error & $1.48\times10^{-6}$ & Numerical convergence \\
\bottomrule\end{tabularx}
\end{table}

\section{Noisy landscapes and the stability gap}
The stability theorem for persistence diagrams is valuable because a small uniform perturbation of a scalar function constrains the bottleneck distance between diagrams. But the current use case introduces multiple uncertainties at once: observed counts are noisy; a low-dimensional embedding may distort distances; the true drift may be non-gradient; the potential may be misspecified; and the minimum-saddle pairing relevant to a transition can change under perturbation. The theory does not erase these distinctions.

We study one controlled perturbation: add Gaussian observation noise to sampled values of the synthetic quartic potential, recompute the grid-based persistence, and assess whether the expected source--saddle feature is identified. At zero noise, the median estimated persistence is 0.24994 and the target pairing is found in 100\% of replicates. At standard deviation 0.01, median persistence is 0.25910 and pairing rate 72.5\%. At noise 0.03, the median is 0.28643 and pairing rate 35\%; at noise 0.12, median persistence reaches 0.42190 and pairing rate is only 32.5\%. The median absolute error increases from about $5.95\times10^{-5}$ without noise to 0.17190 at noise 0.12.

This is a warning, not a universal estimate of the noise robustness of persistent homology. The simple grid and pairing procedure is intentionally transparent but crude. A production method should use a stated smoothing estimator, confidence sets for critical points, bootstrap persistence, sensitivity to embedding and metric, and comparisons across plausible landscapes. The numerical trend cannot be explained away by a stability theorem because the relevant function error, sampling scheme, and feature correspondence are not controlled as in the theorem's hypotheses.

\fig{0.84}{fig03_persistence_robustness}{Synthetic perturbation study with 40 replicates per noise level. As observation noise increases, the estimated persistence moves upward and the intended minimum--saddle pairing becomes less frequently identifiable. Error quantiles and pairing rates are outputs of the released script.}{fig:persistenceNoise}

\subsection{Implication for control}
A protocol optimizer should not treat a persistence estimate as a precisely known barrier. One defensible design is a robust objective over an ensemble of fitted landscapes, with a penalty for policies that succeed only under a fragile persistence pairing. Another is to require the relevant interval to be stable over a prespecified range of preprocessing choices. Both choices trade performance for epistemic robustness and should be reported explicitly. The current study does neither policy-level uncertainty propagation nor real-data bootstrapping, so that stage remains future work.

\section{Non-gradient drift and Hodge decomposition}
A gradient system is only one special case. In a domain with an appropriate inner product and boundary conditions, a vector field can be decomposed into gradient, co-exact/curl, and harmonic components. In a schematic Euclidean notation,
\begin{equation}
b(x)=-\grad V(x)+c(x)+h(x),
\label{eq:hodge}
\end{equation}
where $c$ is a rotational component orthogonal to gradients under the chosen Hodge inner product and $h$ is a harmonic component associated with global topology and boundary conditions. This schematic equation hides important functional-analytic details: the decomposition is not unique until the domain, metric, boundary conditions, and function spaces are specified. A discrete state-transition graph has its own edge-flow decomposition, which is not automatically equivalent to a continuous gene-expression vector field.

The central interpretive caution is that local circulation and global topology are different objects. A planar vector field can have nonzero local curl even on a simply connected domain. A space with nontrivial $H_1$ can host loops, yet a specific endpoint-control problem on it can still have a constant optimal input. Thus the testable hypothesis is not “curl creates pulses” but rather: under a stated control-affine dynamics, dose constraint, noise model, endpoint requirement, and set of competing schedules, a measurable curl-related statistic predicts the relative value of temporal variation beyond information already contained in drift, noise, and state representation.

\fig{0.78}{fig08_toggle_network}{Schematic toggle-switch regulatory network used to orient the synthetic model class. The graph is conceptual and does not encode experimentally inferred edge strengths.}{fig:toggle}

\fig{0.78}{fig09_cycle_network}{Directed cycle graph illustrating circulation as a structural motif. This graph does not establish that the neural-state landscape has a corresponding nontrivial harmonic component; that requires a defined complex, metric, and data-supported flow.}{fig:cycle}

\subsection{A circle counterexample}
Take the state space $S^1$ with angular coordinate $\theta$, natural drift $\omega$, and scalar control $u$:
\begin{equation}
\dot\theta(t)=\omega+u(t),\qquad \theta(0)=0,\quad \theta(T)=\Theta.
\end{equation}
Among square-integrable controls satisfying the endpoint condition, minimize $E(u)=\int_0^T u(t)^2\,\D t$. Since the endpoint imposes $\int_0^T u(t)\D t=\Theta-\omega T$, Cauchy--Schwarz gives
\begin{equation}
E(u)\geq\frac{(\Theta-\omega T)^2}{T},
\end{equation}
with equality if and only if $u$ is constant almost everywhere. For the declared parameters $T=3$, $\omega=0.2$, and $\Theta=1$, the optimal constant is $u^*=0.13333$ and its energy is 0.05333. A rectangular pulse active for 30\% of the horizon and chosen to satisfy the same endpoint has energy 0.17778, a ratio of 3.333. The space has first Betti number $b_1=1$, yet the constant input wins exactly.

\fig{0.84}{fig05_rotor_counterexample}{Analytic counterexample on $S^1$: a nontrivial loop does not force a pulsed control. With endpoints matched, the constant control minimizes quadratic energy and the 30\%-duration pulse costs 3.33 times as much.}{fig:rotor}

This counterexample falsifies an unconditional claim that any H\textsubscript{1} cycle between source and target necessarily makes constant dosing suboptimal. It does not rule out pulsing in non-gradient multi-attractor systems, state-dependent input channels, thresholded success objectives, or environments where a pulse crosses a transient controllability window. Those are narrower conditional questions that must be tested case by case.

\section{Schrödinger bridges as a population-level objective}
Single-cell sequencing generally samples different cells at different time points. A population snapshot provides an empirical distribution, not a movie of the same cell. If $\mu_0$ and $\nu_T$ are source and target state distributions and $\mathbb{P}$ is a reference path measure induced by an uncontrolled diffusion, the Schrödinger bridge selects a path measure $\mathbb{Q}$ that minimizes relative entropy subject to endpoint marginals:
\begin{equation}
\mathbb{Q}^{\star}=\argmin_{\mathbb{Q}:\,\mathbb{Q}_0=\mu_0,\,\mathbb{Q}_T=\nu_T} D_{\mathrm{KL}}(\mathbb{Q}\Vert\mathbb{P}).
\label{eq:sb}
\end{equation}
Under compatible diffusion assumptions, this path-space KL objective corresponds to a stochastic control effort. The exact conversion depends on diffusion and normalization; a model cannot use the phrase “minimum energy” without specifying the reference process and control metric. Soft endpoint penalties, finite samples, and misspecified noise change the optimization problem.

For cell-state control the bridge is a useful complement to an individual-target hitting probability. It asks for a low-relative-entropy evolution of an entire population while matching target marginals. This can discourage an optimizer from sending only a small fraction of cells to the desired state while ignoring the rest. On the other hand, a matching bridge between snapshots does not by itself identify cell lineages or establish a causal pathway. Multiple path distributions can share observed marginals, and the answer depends on the reference dynamics and cost.

A plausible topology-aware objective could add a robust penalty for uncertain basin escape and a schedule-regularity penalty, for example
\begin{equation}
J(u)=\lambda_{\mathrm{SB}}D_{\mathrm{KL}}(\mathbb{Q}^{u}\Vert\mathbb{P})+\lambda_D D(u)+\lambda_S\mathcal{R}(u)+\lambda_U\mathcal{U}_{\mathrm{landscape}}(u),
\label{eq:objective}
\end{equation}
where $\mathcal{R}$ measures temporal complexity and $\mathcal{U}_{\mathrm{landscape}}$ is a penalty for sensitivity to uncertainty over fitted dynamics or persistence pairings. This is a design template rather than a solved theorem. The weights must be chosen and sensitivity-analyzed; the objective could also create perverse trade-offs if target mass, toxicity, and uncertainty are not defined in measurable terms.

\begin{table}[tbp]\centering\small
\caption{Objective terms and the specific question each term answers.}
\label{tab:objective}
\begin{tabularx}{\linewidth}{@{}p{0.19\linewidth}p{0.39\linewidth}X@{}}\toprule
Term & Role & Limitation \\\midrule
Bridge KL & Keeps controlled paths close to reference dynamics while matching marginals & Depends on reference process and diffusion geometry \\
Dose $D(u)$ & Limits total intervention exposure & A scalar budget is not a toxicity model \\
Schedule regularity $\mathcal{R}$ & Penalizes rapid switching or complex pulse trains & Regularization can rule out a genuinely useful transient input \\
Landscape uncertainty & Discourages schedules that depend on unstable inferred basins & Requires a calibrated ensemble; not implemented in the toy RL benchmark \\
Terminal distribution loss & Rewards population-level state matching & Marginals may not identify unique lineages or mechanisms \\
\bottomrule\end{tabularx}
\end{table}

\section{Neural stochastic dynamics: what was fitted}
A data-driven drift model can support control only if its predictive accuracy is evaluated on transitions not used for fitting, noise is estimated separately, and errors in state regions important to the control policy are characterized. We generated a synthetic stochastic dataset from the declared simulator, split 7,200 transitions for training and 1,800 for testing, then fitted an MLP regression model to predict drift. The test drift RMSE is 0.01169, and test $R^2$ is 0.999927. The estimated diffusion standard deviation is 0.08146 against a synthetic truth of 0.08.

\fig{0.88}{fig04_neural_drift_fit}{Synthetic neural drift fit. The held-out transitions are generated by the same known simulator family as training data, so the high score tests implementation and interpolation under a favorable synthetic distribution; it is not evidence of generalization to experimental gene-expression data.}{fig:driftfit}

These numbers are useful as software checks but are deliberately not framed as biological performance. The train/test split is drawn from a common simulator rather than a held-out perturbation family, different cell line, donor, or laboratory. A meaningful biological test would require independent perturbation settings and donors; calibration of process versus measurement noise; uncertainty-aware prediction intervals; and comparison against simple baselines such as linear drift, generalized additive models, or mechanistic gene-regulatory equations. Predictive fit must be evaluated near the boundary of the target basin, since average RMSE can hide errors where control decisions are most sensitive.

\subsection{Minimal fit protocol for real data}
For a future single-cell application, the analysis should start from an explicitly documented processed matrix and sample metadata, preserve donor and batch groupings, and pre-register the state representation. It should avoid leakage from future time points into drift training. Model choice should be compared with linear and mechanistic baselines. Residual diffusion should not be assumed isotropic unless diagnostics support that simplification. The fit should be accompanied by posterior or bootstrap ensembles, calibration metrics, out-of-distribution tests, and an explicit stop rule if the predicted drift is not identifiable from cross-sectional measurements.

\section{Synthetic controller benchmark}
The primary control test uses an abstract noisy bistable state $x_t$ and an intervention with a matched nominal dose upper bound of 3.2. We compare four schedules: constant dosing, evenly spaced pulses, a myopic greedy-pulse controller, and a tabular Q-learning policy. Each benchmark contains 400 evaluation episodes using declared seeds and the same target condition. The pulse schedule uses up to eight pulses, and the constant schedule distributes the same nominal capacity over the full horizon. The objective is endpoint success at the end of the simulation horizon, not biological reprogramming efficiency.

The key result is negative for the RL hypothesis. Evenly spaced pulses succeed in 93/400 episodes (23.25\%), compared with 31/400 (7.75\%) for the trained tabular RL policy, 11/400 (2.75\%) for the greedy policy, and 0/400 for constant dosing in this one parameter regime. The RL policy also uses a mean nominal dose of 3.073 against the budget ceiling 3.2, whereas the even pulse policy uses 3.165. Because pulse use and constant dose are constrained in a simplified simulator, the empirical ordering is a property of these chosen dynamics and reward function; it is not a general conclusion about real cell protocols.

\begin{table}[tbp]\centering\small
\caption{Matched-budget synthetic policy benchmark. All values are from 400 simulator rollouts per policy; no cell data are involved.}
\label{tab:policyresults}
\begin{tabularx}{\linewidth}{@{}Xrrrr@{}}\toprule
Policy & Successes & Success rate & Mean dose & Mean terminal $x$ \\\midrule
Constant-dose & 0/400 & 0.000 & 3.200 & $-0.6672$ \\
Evenly spaced pulses & 93/400 & 0.2325 & 3.165 & 0.2469 \\
Greedy pulses & 11/400 & 0.0275 & 3.200 & $-0.8219$ \\
Tabular RL policy & 31/400 & 0.0775 & 3.073 & 0.0045 \\
\bottomrule\end{tabularx}
\end{table}

\fig{0.9}{fig06_policy_trajectories}{Median synthetic state trajectories under the evaluated policies. Endpoint and state paths come from the simulator; they should not be interpreted as gene-expression trajectories.}{fig:trajectories}

\fig{0.86}{fig07_rl_learning_curve}{Training-return record for the tabular RL agent. The displayed learning curve is diagnostic of the implemented synthetic objective, not proof of convergence to a global optimum. The simple pulse heuristic remains a stronger test-set baseline.}{fig:learning}

\subsection{Interpretation of the controller failure}
The current result is not rescued by pointing to the RL policy's mean terminal state being closer to zero than the pulse policy. The success indicator is defined by target-set membership, so the metric relevant to the benchmark is the success probability, with dose and terminal state as secondary diagnostics. The policy was trained under a reward that mixes success, failure, progress, and exposure penalty. Different reward shaping or training budgets might change its ranking, but tuning until RL wins would convert a falsifiable comparison into a desired outcome. The published result should preserve this failed comparison and motivate stronger baseline coverage.

The evenly spaced schedule is a heuristic, not an oracle. Before making claims about policy quality, future tests should include dynamic programming on a discretized known-state model, model-predictive control with a validated dynamics model, cross-entropy method, stochastic optimal control where tractable, and multiple independent RL seeds. The policy comparison also needs uncertainty intervals and a pre-registered primary outcome. Here the observed binomial standard errors are approximately 0.021 for the pulse schedule and 0.013 for RL; even these simple uncertainty measures show that the current finite simulation is not precise enough to rank small differences.

\section{Dose matching and schedule sensitivity}
Equal nominal dose is an important but incomplete control. The constant schedule uses a low level at every step. A pulse schedule can use high amplitude for fewer steps, and the same L\textsubscript{1} integral does not mean the same internal state exposure when the system is nonlinear. Nor does it control the time derivative, frequency content, duration above threshold, or toxicity. We therefore define the dose budget explicitly, keep the horizon fixed, and report the amplitude bound used by the pulse baseline. This is a fairer comparison than comparing protocols with unequal total administered input, but it is not a complete experimental equivalence.

The sensitivity grid spans five values each for the adaptation parameter and noise scale, with 40 trials per policy and cell (50 cells in total). When the adaptation parameter is zero, both simple policies succeed nearly always at the lowest-noise settings. For adaptation parameter 6—the setting of the headline benchmark—the evenly spaced pulses are stronger than constant dosing across the reported noise grid; by noise standard deviation 0.12, their success frequency is 0.35 versus 0 for constant dosing. The result changes with parameter settings, so no general claim that temporal variation always wins is warranted.

\fig{0.88}{fig12_schedule_sensitivity}{Success-frequency sensitivity sweep over the synthetic adaptation and noise grid. Every cell has 40 trials. This plot reveals parameter dependence rather than a universal advantage for pulsing.}{fig:sensitivity}

\begin{table}[tbp]\centering\small
\caption{Selected sensitivity outcomes at adaptation parameter $\beta=6$. Each percentage is a frequency from 40 simulated trials.}
\label{tab:sensitivity}
\begin{tabularx}{\linewidth}{@{}rrrX@{}}\toprule
Noise SD & Constant success & Pulse success & Interpretation \\\midrule
0.02 & 0/40 & 0/40 & Difficult target in this regime \\
0.05 & 0/40 & 3/40 & Pulse advantage remains low in absolute terms \\
0.08 & 0/40 & 6/40 & Consistent with main-benchmark ordering \\
0.12 & 0/40 & 14/40 & Largest pulse frequency in the selected range \\
0.16 & 1/40 & 18/40 & Success rises with noise for this model; not a biological law \\
\bottomrule\end{tabularx}
\end{table}

The increase in success frequency with noise at high adaptation is a model-dependent stochastic-resonance-like pattern, not a validated biological claim. It should be compared against alternative reward definitions and trajectory metrics. A basic check is whether apparent success is driven by broad noise-induced target crossing rather than control. That comparison is not complete here and is an explicit limitation.

\section{Evaluation protocol and baselines}
We separate three layers of evidence: (i) analytic identities and counterexamples, (ii) synthetic model tests, and (iii) real-data and ultimately experimental tests. In the current paper only the first two have results. The public single-cell resource is identified for the future third layer, but its matrix is not downloaded or analyzed in the current cycle. This separation is central to the interpretation of every metric.

\fig{0.93}{fig10_pipeline_flow}{Workflow diagram from model specification through synthetic verification and a separately gated future public-data analysis. The public-data branch is a plan only; it is not shown as completed.}{fig:pipeline}

A robust evaluation should include: a gradient potential with known barrier; a non-gradient system with a known quasipotential or numerical action solution; a circular or graph-structured state space; perturbation experiments over data noise; multiple control budgets; a policy test with matched action constraints; and a counterexample to each proposed implication. We implemented only a subset of this matrix. The completed analytic tests include the quartic landscape and circle endpoint problem; the simulation covers noisy persistence, synthetic drift fitting, and one policy regime plus a sensitivity grid. Full multi-attractor networks, uncertainty propagation from neural drift into policies, and a genuine Schrödinger bridge solver are not implemented.

\begin{table}[tbp]\centering\small
\caption{Validation inventory with explicit status.}
\label{tab:validation}
\begin{tabularx}{\linewidth}{@{}p{0.31\linewidth}p{0.20\linewidth}X@{}}\toprule
Test & Status & Evidence and limitation \\\midrule
Quartic gradient barrier vs H\textsubscript{0} & PASS & Analytic barrier and numerical pairing agree to $1.69\times10^{-6}$ \\
FW action discretization convergence & PASS & Error decreases over grid refinement in the toy model \\
Persistence noise robustness & PASS as executed; robustness fails & 240 replicates reveal pairing and calibration deterioration \\
Synthetic neural drift fit & PASS as executed & Held-out synthetic transitions only; no empirical generalization \\
Constant vs pulse schedules & PASS as executed; RL hypothesis not supported & Even pulse policy beats RL under selected benchmark \\
Loop implies pulses & FALSIFIED in general & Exact circle example has constant-energy optimum \\
Public single-cell reprogramming fit & NOT RUN & Dataset identified; data not downloaded or processed \\
Schrödinger bridge controller & NOT IMPLEMENTED & Objective formalized; no solver run \\
Wet-lab testing & NOT RUN / outside scope & No biological efficacy claim or protocol recommendation \\
\bottomrule\end{tabularx}
\end{table}

\section{Falsifiable predictions and decisive negative tests}
The paper proposes a restrained set of predictions that would make a future study scientifically useful even if the framework fails. They are scoped to models and datasets, not phrased as clinical or wet-lab instructions.

\subsection{Prediction A: barrier--persistence agreement has a conditional domain}
In systems known to be close to gradient diffusions, and when the source minimum is correctly paired with the saddle separating it from the target basin, bootstrap persistence estimates should track an independently estimated quasipotential barrier within a tolerance that accounts for potential-fit and discretization error. The prediction is falsified if the relationship fails after controlling state representation, metric, sample size, and model uncertainty—or if a non-gradient alternative consistently predicts transition rates better from dynamical action than from persistence. This test must compare at least one non-gradient system, not only repeat the same quartic toy example.

\subsection{Prediction B: schedule benefit requires dynamic evidence beyond Betti number}
A cycle-related measure should predict out-of-sample benefit from temporal variation only after controlling for state-space dimension, local drift magnitude, controllability Gramian, target distance, and noise covariance. The circle counterexample already falsifies the unconditional version. A test that uses $b_1$ alone as a predictor should be considered failed if it does not outperform random topology-label permutations or if local curl explains the result equally well in simply connected domains.

\subsection{Prediction C: RL must generalize beyond the training simulator}
If RL provides value, its learned policy must improve over equally budgeted pulse, greedy, model-predictive, and dynamic-programming baselines on held-out parameter regimes, random seeds, and model classes. Reporting only training returns would be insufficient. A fair comparison freezes policy hyperparameters before test evaluation, uses confidence intervals, and reports all attempted seeds. If the RL policy's mean performance is lower than a simple heuristic across the prespecified test grid, the RL-superiority claim is rejected even if one tuned setting shows a win.

\subsection{Prediction D: uncertainty-aware topology helps, or it should be dropped}
Landscape uncertainty is useful only if carrying an ensemble through policy optimization improves calibration or robustness against model misspecification. If it adds computational complexity without improving held-out success or reducing worst-case failure, it should not remain in the objective merely to preserve theoretical elegance.

\fig{0.91}{fig11_falsification_flow}{Falsification path: validate the scalar/dynamical correspondence, test the cycle hypothesis against counterexamples, compare controllers with a prespecified benchmark, then move to public single-cell data only if the synthetic gates survive.}{fig:falsify}

\section{Failure modes, robustness, and uncertainty}
The most immediate issue is landscape non-identifiability. A scalar potential inferred from stationary density assumes a relationship between drift and noise. In a non-equilibrium steady state, probability currents can exist, and the stationary density alone does not identify those currents. A learned energy-like scalar can fit marginal density while failing to predict transition paths. Consequently, a persistence score should be treated as a feature of the estimation pipeline—not as a directly observed molecular energy.

Second, graph topology depends on how the data are transformed. Principal-component projections, diffusion maps, and latent models can merge distinct cell states or introduce apparent loops. Feature filtration and nearest-neighbor parameters also change the persistence diagram. Stability of the persistence diagram under perturbations to a fixed scalar function is not a substitute for uncertainty over the embedding and the learned drift. In the synthetic noise test, pairing rates deteriorate much faster than one might infer from an eyeballed smooth landscape.

Third, the controller's objective may be a poor proxy for actual success. Endpoint-only reward can ignore transient loss of viability, undesired intermediate states, off-target differentiation, and whether a trajectory is reversible. A nominal dose budget does not equal exposure cost, and the target basin in a toy scalar model has no direct biological interpretation. A future constrained objective should include feasibility and safety-related filters informed by experimental teams, not a reward function optimized in isolation.

Fourth, RL can exploit simulator artifacts, reward shaping, and distribution shifts. The present RL policy did not beat the evenly spaced pulse policy, which is precisely why both should remain in the evidence record. More iterations do not automatically fix model misspecification. A policy that wins only when the transition model is exact can fail when the biological system differs from the training simulator. Robust evaluation must vary drift parameters, diffusion, response delay, and action nonlinearities before a policy earns any real-data claim.

\begin{table}[tbp]\centering\small
\caption{Failure modes and required mitigations.}
\label{tab:failures}
\begin{tabularx}{\linewidth}{@{}p{0.22\linewidth}p{0.34\linewidth}X@{}}\toprule
Failure mode & How it misleads & Minimum mitigation \\\midrule
False potential interpretation & Scalar score is treated as physical energy & Fit and compare dynamical models; test current and non-gradient alternatives \\
Incorrect persistence pairing & Interval does not correspond to source--target saddle & Track critical points, paths, and pairing across uncertainty ensembles \\
Embedding-induced topology & Projection creates/removes apparent loops & Compare embeddings and metrics; test against known synthetic ground truth \\
Noise-misspecified SDE & Escape action is estimated in wrong metric & Estimate diffusion covariance and validate residuals; propagate uncertainty \\
RL simulator exploitation & Policy succeeds on training dynamics only & Frozen held-out regimes, strong baselines, multiple seeds, adversarial models \\
Nominal dose mismatch & Equal integral hides unequal biological exposure & Report amplitude, duration, regularity, and toxicity proxies separately \\
Endpoint reward pathology & Intermediate failures are ignored & Add trajectory constraints and report intermediate state occupancy \\
Cross-sectional ambiguity & Distribution transport mistaken for lineage evidence & State explicitly that population marginals do not identify unique trajectories \\
\bottomrule\end{tabularx}
\end{table}

\section{A public-data path that is not yet executed}
A relevant public resource is GEO GSE242423, which documents single-cell RNA-sequencing across a 14-day OSKM-induced human-fibroblast reprogramming time course with samples labelled D0, D2, D4, D6, D8, D10, D12, D14, and iPSC. The associated super-series GSE242424 includes related chromatin-accessibility samples and multiome measurements. GEO's record notes that raw data are not provided there due to privacy restrictions, while processed data are listed as supplementary files \citep{geoGSE242423,geoGSE242424}. Identifying the accession gives the project a plausible data target; it does not mean that the data were fetched, processed, or legally/technically accessible in this research run.

The next empirical-oriented phase should first verify the accessions, license/usage conditions, sample metadata, donor/batch structure, feature matrices, and exact availability of processed data. It should then predefine the latent state representation and the interpretation of the target distribution. Since cells are destructively sampled, transitions must be inferred as population-level stochastic evolution, not presented as observed single-cell lineages. If processed counts and metadata can be acquired, train drift models only on allowed time slices, validate on held-out times or experiments, and compare the bridge predictions against standard OT and a mechanistic bursty-reference model. The 2026 study by Fournié and colleagues provides a particularly important comparator because it formulates a Schrödinger problem with a mechanistic stochastic gene-expression reference \citep{fournie2026schrodinger}.

A public-data analysis should be abandoned or narrowed if access, batch structure, measurement timing, or model identifiability do not support the desired estimand. In particular, time-stamped snapshots alone may not identify a unique drift field or noise covariance. The right output in that case is a carefully scoped identifiability statement, not a forced policy benchmark. No new biological observations are reported in this preprint.

\section{Synthetic network models and model-class extension}
Toggle-switch networks offer a compact way to create multiple attractors: mutually repressive regulators can support alternative stable expression states, depending on parameters and molecular noise. Ring-like repressilator networks illustrate oscillatory dynamics rather than a scalar basin picture. These architectures are familiar examples, not newly proposed biological mechanisms. They are useful because they expose the difference between basin topology, local rotational dynamics, and temporal protocol response.

The current scripts draw network schematics as conceptual diagrams and do not simulate those network graphs as the policy benchmark. The paper therefore does not claim that the controller result generalizes from the scalar bistable SDE to a toggle switch or repressilator. A credible extension would implement known synthetic regulatory equations, compare deterministic attractors with stochastic transitions, calculate graph- or state-space topology from sampled trajectories, and test whether a topology-derived metric adds predictive value beyond drift and controllability features. It would also compare model-generated observations with the true simulator state to quantify the effect of measurement noise.

\subsection{A minimum suite of synthetic ground truths}
We propose four escalating classes. Class I is the quartic one-dimensional gradient potential, used for analytical calibration. Class II is a non-gradient rotor or planar system with known circulation, used to test the cycle hypothesis. Class III is a multi-attractor regulatory network with known governing parameters, used to distinguish true basins from projection artifacts. Class IV is a partially observed stochastic system with unknown parameters, used for model-identification stress tests. A success in Class I is necessary for implementation verification but nowhere near sufficient for biological utility.

\section{Control-theoretic conditions for schedules}
The optimal control law depends not only on topology but on the control channel $B(x)$, the endpoint set, time horizon, amplitude limits, state-dependent drift, noise covariance, and objective. If the system is linear with quadratic cost and unconstrained endpoints, the optimal controller may be smooth and continuous. If input is bounded, response delay matters, or the system has a threshold and strongly nonlinear drift, bang-bang or pulsed controls may arise. Those features are control-theoretic, not consequences of the first Betti number by itself.

An appropriate formal test would compare the value functions for constant and time-varying controls:
\begin{equation}
J_{\mathrm{const}}^*=\inf_{u_t\equiv u_0}J(u),\qquad
J_{\mathrm{time}}^*=\inf_{u\in\mathcal{U}_{\mathrm{ad}}}J(u).
\end{equation}
Temporal variation offers a strict advantage only if $J_{\mathrm{time}}^*<J_{\mathrm{const}}^*$ for the stated cost and constraints. No topological summary can replace that inequality. To show topology is a useful predictor, one would need an identified statistic $\tau(\mathcal{X},b,\sigma,B)$ and a tested relation between $\tau$ and the gap $J_{\mathrm{const}}^*-J_{\mathrm{time}}^*$ across a prespecified collection of systems.

A pulse-count lower bound is even more demanding. Pulse count is not well-defined until amplitude thresholds, pulse segmentation, minimum gaps, and allowable bandwidth are specified. A time-varying function can be approximated by infinitely many control parameterizations. Any lower bound must therefore include regularity constraints, a model of input kinetics, and a precise admissible control class. A topological integer cannot by itself determine a number of pulses unless a mathematical mapping from topology to a necessary control decomposition has been proved. We do not prove such a mapping, and the current evidence does not support it.

\section{Discussion: what the evidence allows us to say}
The conceptual appeal of connecting persistent homology, Hodge decomposition, stochastic bridges, and RL is that each describes a different aspect of the problem: connectivity across energy thresholds, rotational versus potential-like flow, population-level path distributions, and sequential policy search. The difficulty is exactly that these quantities live at different mathematical levels. Joining them produces an explanatory hypothesis, not a result. A strong paper has to test whether the connection adds information instead of simply placing the terms beside one another.

The quartic potential demonstrates a real but narrow identity: with a valid scalar potential and correct pairing, H\textsubscript{0} persistence equals the barrier height. The noise sweep demonstrates fragility of the simplest estimator under observation perturbation. The circle example decisively rules out the broad loop-implies-pulsing claim. The controller benchmark adds a second negative result: under this model and reward, RL fails to beat a straightforward pulse heuristic. Together, these findings argue for a pipeline that uses topology as an uncertain diagnostic, retains explicit dynamical information, and keeps simple baselines central.

The remaining research opportunity is still meaningful. A carefully identified Hodge-like feature might complement controllability and path action when predicting schedule benefit. A bridge objective may help match full populations better than endpoint hit-rate optimization. And RL might be valuable in a richer environment where policies need to adapt to state estimates and uncertain response. Each claim should remain a hypothesis until it beats alternatives in held-out models and eventually in suitable experimental datasets.

\subsection{Human benefit and translation boundary}
If future work identifies a robust relationship between state-space structure and intervention schedule, it could help experimentalists prioritize candidate reprogramming schedules and reduce uninformative trial-and-error. The causal pathway is indirect: reliable state inference would inform an in-silico schedule, which would then need independent experimental testing and safety assessment. The present work is far from those stages. It cannot specify doses, predict clinical efficacy, guarantee cell quality, or justify a protocol in human cells. There are significant barriers in measurement noise, batch and donor heterogeneity, biological safety, off-target states, ethics, and reproducibility.

\section{Limitations}
The principal limitation is scope of validation. All numerical outputs are synthetic and driven by known model equations. No public single-cell expression data have been downloaded or fitted. No wet-lab experiments, cellular assays, or clinical outcomes are included. Consequently, the manuscript does not show that the proposed combination works on real cellular reprogramming.

Second, the neural drift estimator is trained and tested on simulated transitions generated from the same model family. Its excellent test score is an implementation benchmark under favorable assumptions. It does not demonstrate the identifiability of drift from sparse single-cell snapshots, nor does it model bursty transcription, nonstationary noise, donor effects, or measurement artifacts.

Third, the topology code uses a discretized one-dimensional potential and a simple pairing procedure. While adequate for a known quartic function, it does not solve the challenge of high-dimensional persistent homology on latent embeddings. The robustness experiment adds noise to a one-dimensional function and reveals poor feature correspondence; it is illustrative, not a calibrated error model for RNA sequencing.

Fourth, the Hodge decomposition is conceptual, and the scripts do not estimate a Hodge field from observed trajectories. Neither the Hodge curl nor first homology is connected empirically to the chosen control simulator. The circle example is a mathematical counterexample, not a cellular model.

Fifth, the bridge optimization is formulated but no numerical Schrödinger bridge is solved. The RL implementation uses a tabular agent and a scalar state. The finite policy benchmark has no confidence intervals from independent training seeds, and the policy was not tuned against all strong baselines. These limitations narrow the strength of the conclusions but do not justify hiding the negative outcome.

Finally, the literature search was focused and current as of 24 October 2026 but not exhaustive across every database, patent collection, or unpublished result. Novelty is therefore deliberately stated as an integration and stress-test contribution, not a priority claim.

\section{Conclusion}
Topological summaries can constrain interpretations of a scalar landscape, but their relationship to control cost is conditional. In a gradient diffusion with a correctly paired source minimum and saddle, barrier height and H\textsubscript{0} persistence are the same height difference. In noisy inferred landscapes, the correct pairing and calibration can deteriorate. In non-gradient settings, a scalar persistence value need not equal the quasipotential barrier. A nontrivial loop does not require a pulsed optimum, as the analytic circle counterexample shows. In the synthetic controller benchmark, RL does not beat evenly spaced pulses, and the paper reports that failure rather than claiming success.

The defensible next step is a preregistered, uncertainty-aware comparison across synthetic ground-truth networks followed—only if the synthetic gates survive—by a reproducible analysis of public time-course data. The project's useful hypothesis is not that topology dictates a protocol, but that topological features may become informative control features under explicit, testable dynamical assumptions. Until those assumptions and out-of-sample advantages are demonstrated, topology, Hodge structure, bridges, and RL remain complementary tools, not a proven recipe for cellular reprogramming.

\medskip\noindent\textbf{Co-equal research-process credit is given to Devanik Debnath and Xylia, the autonomous multi-model research engine coordinating AI-assisted discovery, synthesis, and verification. Devanik Debnath retains responsibility for the final manuscript, scientific claims, citations, and released artifacts.}

\clearpage
\appendix
\section{Derivation of the gradient action identity}
For the gradient SDE in Eq.~\eqref{eq:gradientSDE}, the normalized action is
\begin{equation}
S_T(\phi)=\frac14\int_0^T\|\dot\phi+\grad V(\phi)\|^2\D t.
\end{equation}
The pointwise identity
\begin{equation}
\|a+b\|^2=\|a-b\|^2+4a\cdot b
\end{equation}
with $a=\dot\phi$ and $b=\grad V(\phi)$ yields
\begin{equation}
S_T(\phi)=\frac14\int_0^T\|\dot\phi-\grad V(\phi)\|^2\D t+\int_0^T\dot\phi\cdot\grad V(\phi)\D t.
\end{equation}
By the chain rule, the second term is $V(\phi(T))-V(\phi(0))$. If the path starts at $x_m$ and ends at $x_s$, the nonnegative squared term gives $S_T\geq V(x_s)-V(x_m)$. Equality occurs when $\dot\phi=\grad V(\phi)$, the uphill time reversal of the deterministic gradient flow $\dot x=-\grad V(x)$. Strict equality over a finite interval depends on endpoint asymptotics and whether the minimizing path is admissible; the action infimum still agrees with the barrier under the usual idealized assumptions.

This derivation also illustrates why the statement is not universal. If the drift is not $-\grad V$, the cross-term does not reduce to the differential of a scalar potential. If the noise covariance is $a(x)=\sigma(x)\sigma(x)^\top$, the action is weighted by the inverse covariance on its range. Such a metric-weighted action cannot be compared with scalar persistence without further structure.

\section{Derivation of the circular-control counterexample}
With $\dot\theta=\omega+u$, endpoint feasibility imposes $\int_0^T u(t)\D t=\Theta-\omega T=:c$. By Cauchy--Schwarz,
\begin{equation}
\left(\int_0^T u(t)\D t\right)^2\leq T\int_0^T u(t)^2\D t,
\end{equation}
so $E(u)\geq c^2/T$. Equality requires $u(t)=c/T$ almost everywhere. For $T=3$, $\omega=0.2$, and $\Theta=1$, $c=0.4$, and the minimum energy is $0.4^2/3=0.053333\ldots$. If a pulse acts for a fraction $p=0.3$ of the horizon, endpoint feasibility requires pulse amplitude $c/(pT)=0.4/0.9=0.4444\ldots$; its energy is $pT[c/(pT)]^2=c^2/(pT)=0.177777\ldots$. The ratio to the constant-control energy is $1/p=10/3$. Since $S^1$ has $b_1=1$, this exact construction is a counterexample to a topology-only necessary-pulse claim.

\clearpage
\section{Synthetic model and benchmark specification}
The release code is the authoritative executable specification for simulator parameters, random seeds, action bounds, reward function, rollout logic, training settings, and plot generation. The model is intentionally small enough to run deterministically without downloading any data. The potential is $V(x)=(x^2-1)^2/4$. A drift estimator is trained on 7,200 synthetic transitions and tested on 1,800 transitions; the policy study uses 400 evaluation rollouts for each of four controllers. Sensitivity results contain 50 parameter--noise--policy cells with 40 trials per cell. The random seed for the main verification is 20261024; the neural drift training subroutine uses seed 20261035 as recorded in the machine-readable report.

The verifier checks the analytic barrier, persistence, action convergence, neural-fit metrics, dose bounds, circle counterexample, policy outputs, sensitivity grid, replicate retention, and figure generation. Passing certifies code checks only; CSV outputs permit recomputation, and RL superiority over pulses is not asserted.

\section{Minimal public-data analysis protocol}
A public-data follow-up should preserve the following order: (1) verify GSE242423/GSE242424 metadata and access conditions; (2) acquire available processed matrices and checksums; (3) predefine QC, batch handling, and a reduced state representation; (4) split evaluation by time, donor, or experiment to avoid leakage; (5) fit multiple drift/noise models and quantify uncertainty; (6) estimate persistent features across bootstrap samples and embeddings; (7) solve a bridge or transport baseline with documented assumptions; (8) compare constant, simple pulse, model-predictive, dynamic-programming, and RL policies on simulated counterfactuals only where identified; (9) report non-identifiability and negative results; and (10) submit protocol hypotheses for experimental assessment without describing simulated schedules as validated interventions.

The analysis should be stopped if state-space inference or noise estimation is not sufficiently identifiable. The current release has not run any of these data-acquisition and empirical-modeling steps. This appendix is a protocol for a future computational phase, not a claim that the phase has been completed.

\section{Reproducibility and interpretation}
To reproduce the synthetic checks, install the Python dependencies listed in \texttt{results/environment.json} and execute:
\begin{verbatim}
python verify_model.py
\end{verbatim}
The script regenerates output tables and figure assets. Running the script changes generated outputs if a dependency implementation or plotting backend differs; record exact package versions and compare CSV values before interpreting differences. Use a clean working directory as a test of path independence. The PDF build command is documented in the release README and deployment script.

The release deliberately distinguishes three meanings of “validation”: program assertions, agreement with an analytical toy result, and empirical biological confirmation. This work achieves the first two for the named tests and achieves none of the third. The public dataset accession is a planned candidate only. No experiment, cell line, sample, dose, patient, or biological treatment was studied here.

\clearpage
\begin{thebibliography}{99}
\bibitem[Waddington(1957)]{waddington1957strategy}
C.~H. Waddington. \emph{The Strategy of the Genes}. George Allen \& Unwin, London, 1957.

\bibitem[Zhou et~al.(2013)]{zhou2013waddington}
J.~X. Zhou, M.~D.~S. Aliyu, E. Aurell, and S. Huang.
Quantifying Waddington landscapes and paths of non-adiabatic cell fate decisions for differentiation, reprogramming and transdifferentiation.
\emph{Journal of the Royal Society Interface}, 10(89):20130787, 2013. doi: \href{https://doi.org/10.1098/rsif.2013.0787}{10.1098/rsif.2013.0787}.

\bibitem[Edelsbrunner et~al.(2002)]{edelsbrunner2002persistence}
H. Edelsbrunner, D. Letscher, and A. Zomorodian.
Topological persistence and simplification. \emph{Discrete \& Computational Geometry}, 28:511--533, 2002. doi: \href{https://doi.org/10.1007/s00454-002-2885-2}{10.1007/s00454-002-2885-2}.

\bibitem[Cohen-Steiner et~al.(2007)]{cohensteiner2007stability}
D. Cohen-Steiner, H. Edelsbrunner, and J. Harer.
Stability of persistence diagrams. \emph{Discrete \& Computational Geometry}, 37:103--120, 2007. doi: \href{https://doi.org/10.1007/s00454-006-1276-5}{10.1007/s00454-006-1276-5}.

\bibitem[Freidlin and Wentzell(2012)]{freidlin2012random}
M.~I. Freidlin and A.~D. Wentzell. \emph{Random Perturbations of Dynamical Systems}. 3rd edition, Springer, 2012. doi: \href{https://doi.org/10.1007/978-3-642-25847-3}{10.1007/978-3-642-25847-3}.

\bibitem[Kramers(1940)]{kramers1940brownian}
H.~A. Kramers. Brownian motion in a field of force and the diffusion model of chemical reactions. \emph{Physica}, 7(4):284--304, 1940. doi: \href{https://doi.org/10.1016/S0031-8914(40)90098-2}{10.1016/S0031-8914(40)90098-2}.

\bibitem[Jiang et~al.(2011)]{jiang2011hodge}
X. Jiang, L.-H. Lim, Y. Yao, and Y. Ye. Statistical ranking and combinatorial Hodge theory. \emph{Mathematical Programming}, 127:203--244, 2011. doi: \href{https://doi.org/10.1007/s10107-010-0419-x}{10.1007/s10107-010-0419-x}.

\bibitem[Zarzycki et~al.(2025)]{zarzycki2025pbnstac}
J. Zarzycki et al. pbn-STAC: Deep reinforcement learning-based framework for cellular reprogramming. \emph{Theoretical Computer Science}, 1049:115382, 2025. doi: \href{https://doi.org/10.1016/j.tcs.2025.115382}{10.1016/j.tcs.2025.115382}. Code and data: \url{https://github.com/jakub-zarzycki2022/pbn-stac}.

\bibitem[Mizera and Zarzycki(2026)]{mizera2026gattaca}
A. Mizera and J. Zarzycki. The GATTACA framework: Graph neural network-based reinforcement learning for controlling biological networks. \emph{Proceedings of the AAAI Conference on Artificial Intelligence}, 40(2):873--880, 2026. doi: \href{https://doi.org/10.1609/aaai.v40i2.37055}{10.1609/aaai.v40i2.37055}.

\bibitem[Fourni{\'e} et~al.(2026)]{fournie2026schrodinger}
C. Fourni{\'e}, E. Ventre, U. Herbach, A. Baradat, O. Gandrillon, and F. Crauste.
Cell trajectory inference based on schr{\"o}dinger problem and a mechanistic model of stochastic gene expression. \emph{npj Systems Biology and Applications}, 12:120, 2026. doi: \href{https://doi.org/10.1038/s41540-026-00743-x}{10.1038/s41540-026-00743-x}.

\bibitem[Zhang et~al.(2026)]{deciphering2025cellfate}
Z. Zhang, Z. Wang, Y. Sun, J. Shen, Q. Peng, T. Li, et al.
Deciphering cell-fate trajectories using spatiotemporal single-cell transcriptomic data. \emph{npj Systems Biology and Applications}, 12:2, 2026; published online 4 December 2025. doi: \href{https://doi.org/10.1038/s41540-025-00624-9}{10.1038/s41540-025-00624-9}.

\bibitem[NCBI GEO(2023a)]{geoGSE242423}
NCBI Gene Expression Omnibus. GSE242423: Transcription factor stoichiometry, motif affinity and syntax regulate single-cell chromatin dynamics during fibroblast reprogramming to pluripotency [scRNA]. Public accession record, 2023. \url{https://www.ncbi.nlm.nih.gov/geo/query/acc.cgi?acc=GSE242423}. Accessed 24 October 2026. Raw data restrictions are noted in the accession; the current work did not download the data.

\bibitem[NCBI GEO(2023b)]{geoGSE242424}
NCBI Gene Expression Omnibus. GSE242424: Transcription factor stoichiometry, motif affinity and syntax regulate single-cell chromatin dynamics during fibroblast reprogramming to pluripotency. SuperSeries record, 2023. \url{https://www.ncbi.nlm.nih.gov/geo/query/acc.cgi?acc=GSE242424}. Accessed 24 October 2026.

\end{thebibliography}
\end{document}
